\chapter{속성 유형 --- 값의 종류를 먼저 가른다}
\label{ch:attributes}
% source: deck 3 slide 1 (제목만, 내용 없음)

\begin{inbrief}
  \item[Why] 값에 허용된 연산(척도)과 데이터가 놓이는 모양(행렬, 공간)을
             알아야 이후 모든 요약과 그림이 합법이 된다.
  \item[When] 데이터를 받자마자, 어떤 통계도 계산하기 전에.
  \item[Where] 표로 된 모든 데이터. 이미지, 텍스트, 네트워크도 행렬로
               바꾼 뒤 같은 틀에 들어온다.
  \item[How] 열마다 척도와 이산/연속을 가르고, 식별자를 뺀 나머지를
             $n \times d$ 행렬로 본다.
\end{inbrief}

같은 숫자라도 무엇을 뜻하느냐에 따라 허용되는 계산이 다르다. 우편번호
78539와 칼로리 540은 둘 다 숫자로 저장되지만, 우편번호의 평균은 아무 의미가
없다. \term{속성 유형}(Attribute Type)을 먼저 가르는 이유는 여기에 있다.
어떤 요약값이 \emph{합법}인지가 여기서 결정된다.

\begin{keyidea}
속성 유형은 데이터를 예쁘게 분류하려는 작업이 아니다. \textbf{그 값에 어떤
연산이 허용되는지}를 정하는 작업이다. 유형을 틀리면 이후의 평균 $\cdot$ 표준편차
$\cdot$ 그래프가 전부 조용히 틀린다 --- 계산은 성공하고 숫자도 나오기 때문이다.
\end{keyidea}

% =====================================================================
\section{배경}
\label{sec:attr-background}

이 장은 강의 자료 ``Data and Matrix''(덱 3,
\texttt{notes/3\_DataAndMatrix.pptx})를 옮긴 것이다. 책의 첫 본문 장이므로
앞서 다룬 개념은 없지만, 이 장이 바로 쓰는 기초 용어 세 가지는 여기서
한 줄씩 깔아 둔다.

\begin{itemize}
  \item \term{표}(Table): 값을 행과 열의 격자에 늘어놓은 것.
  \item \term{행}(Row)과 \term{열}(Column): 표에서 가로줄이 행, 세로줄이
        열이다.
  \item \term{좌표평면}(Coordinate Plane): 점 하나가 값 두 개(가로 좌표,
        세로 좌표)로 정해지는 평면.
\end{itemize}

이 장은 두 부분으로 나뉜다. 앞부분(\Cref{sec:attr-flow}부터
\Cref{sec:attr-stats}까지)은 열 하나의 성질, 곧 값 하나에 어떤 연산이
허용되는가를 다룬다. 뒷부분(\Cref{sec:attr-discrete}부터)은 그렇게 정리한
열들을 모아 표 전체, 곧 행렬로 보는 법을 다룬다. 이 장에서 정한 유형은
이후 요약값(\Cref{ch:descriptive})과 그래프(\Cref{ch:visualization})를
고르는 전제가 된다.

% =====================================================================
\section{두 번의 질문으로 가른다}
\label{sec:attr-flow}
% source: deck 3 slide 3

\begin{inbrief}
  \item[Why] 값에 허용된 연산을 알아야 요약값과 그래프를 고를 수 있기 때문이다.
  \item[When] 데이터를 받자마자, 어떤 통계도 계산하기 전에 한다.
  \item[Where] 모든 변수에 적용한다. 특히 숫자로 저장된 범주형(우편번호, ID,
               코드)에서 판단이 갈린다.
  \item[How] 질문 두 개를 순서대로 던진다. ``숫자로 계산이 되는가'' $\to$
             ``0이 없음을 뜻하는가'' 또는 ``값들 사이에 순서가 있는가''.
\end{inbrief}

가름은 2단계로 끝난다. 1단계에서 \term{수치형}(Numeric)과
\term{범주형}(Categorical)을 나누고, 2단계에서 각각을 다시 둘로 나눈다
(\Cref{fig:attr-flow}).

\begin{figure}[ht]\centering
\begin{tikzpicture}[
  qbox/.style    = {draw = bookaccent, fill = bookaccentbg, rounded corners = 2pt,
                    minimum width = 46mm, minimum height = 8mm, align = center,
                    font = \small, inner sep = 2pt},
  leafbox/.style = {draw = bookrule, fill = booksurface, rounded corners = 2pt,
                    minimum width = 34mm, minimum height = 11mm, align = center,
                    font = \small, inner sep = 2pt},
  arrow/.style   = {-{Stealth[length=2.2mm]}, draw = bookmuted},
  elab/.style    = {font = \scriptsize\color{bookmuted}, inner sep = 1.5pt}]

  \node[qbox] (q1) at (0, 0) {1단계: 숫자로 계산이 되는가?};

  \node[qbox] (q2) at (-3.9, -1.6) {2단계: 0이 ``없음''인가?};
  \node[qbox] (q3) at ( 3.9, -1.6) {2단계: 순서가 있는가?};

  \node[leafbox] (ratio)    at (-5.8, -3.6) {비율(Ratio)\\[-2pt] \scriptsize 칼로리 $\cdot$ 금액};
  \node[leafbox] (interval) at (-2.0, -3.6) {구간(Interval)\\[-2pt] \scriptsize 섭씨온도 $\cdot$ 연도};
  \node[leafbox] (ordinal)  at ( 2.0, -3.6) {순서(Ordinal)\\[-2pt] \scriptsize 학력 $\cdot$ 만족도};
  \node[leafbox] (nominal)  at ( 5.8, -3.6) {명목(Nominal)\\[-2pt] \scriptsize 혈액형 $\cdot$ 우편번호};

  \draw[arrow] (q1.south) -- ++(0,-0.4) -| (q2.north)
        node[elab, pos=0.45, above] {수치형};
  \draw[arrow] (q1.south) -- ++(0,-0.4) -| (q3.north)
        node[elab, pos=0.45, above] {범주형};

  \draw[arrow] (q2.south) -- ++(0,-0.5) -| (ratio.north)
        node[elab, pos=0.45, above] {0 = 없음};
  \draw[arrow] (q2.south) -- ++(0,-0.5) -| (interval.north)
        node[elab, pos=0.45, above] {0 = 기준점};
  \draw[arrow] (q3.south) -- ++(0,-0.5) -| (ordinal.north)
        node[elab, pos=0.45, above] {서열 있음};
  \draw[arrow] (q3.south) -- ++(0,-0.5) -| (nominal.north)
        node[elab, pos=0.45, above] {서열 없음};
\end{tikzpicture}
\caption{속성 유형을 가르는 두 번의 질문. 1단계에서 수치형/범주형을 나누고,
2단계에서 각각을 다시 둘로 나눈다.}
\label{fig:attr-flow}
\end{figure}

% =====================================================================
\section{수치형 --- 구간과 비율}
\label{sec:attr-numeric}
% source: deck 3 slide 3

\begin{inbrief}
  \item[Why] 같은 수치형이라도 ``몇 배''라는 말이 성립하는 것과 아닌 것이 있기
             때문이다.
  \item[When] 1단계에서 수치형으로 가른 뒤에 던지는 질문이다.
  \item[Where] 온도 $\cdot$ 연도 $\cdot$ 지능지수는 구간, 무게 $\cdot$ 금액
               $\cdot$ 개수는 비율이다.
  \item[How] ``값 0이 정말로 아무것도 없다는 뜻인가''를 묻는다. 그렇다면 비율,
             기준점일 뿐이라면 구간이다.
\end{inbrief}

\begin{definition}[구간 척도와 비율 척도]
\label{def:interval-ratio}
\term{구간 척도}(Interval Scale)는 값의 \emph{간격}은 의미가 있지만 0이
임의의 기준점인 척도다. 덧셈과 뺄셈은 되고, 나눗셈(비율)은 되지 않는다.

\term{비율 척도}(Ratio Scale)는 0이 \term{절대 영점}(Absolute Zero), 즉
``그 속성이 전혀 없음''을 뜻하는 척도다. 덧셈 $\cdot$ 뺄셈에 더해
곱셈 $\cdot$ 나눗셈까지 의미를 가진다.
\end{definition}

\begin{example}[``두 배''가 성립하는가]
\label{ex:twice}
섭씨온도는 구간 척도다. $0^\circ\mathrm{C}$ 는 열이 없다는 뜻이 아니라 물이 어는
지점을 0으로 잡은 약속일 뿐이다. 그래서 $20^\circ\mathrm{C}$ 는
$10^\circ\mathrm{C}$ 보다 ``$10^\circ$ 더 덥다''라고는 말할 수 있어도
``두 배 덥다''라고는 말할 수 없다. 실제로 같은 온도를 화씨로 바꾸면 50\,\textdegree F
와 68\,\textdegree F 가 되어 비율이 $2.0$ 에서 $1.36$ 으로 바뀐다 --- 단위를
바꿨을 뿐인데 답이 달라진다면 그 계산은 애초에 의미가 없었던 것이다.

칼로리는 비율 척도다. 0\,kcal 는 열량이 실제로 없다는 뜻이므로,
540\,kcal 는 270\,kcal 의 ``두 배''가 맞다. 단위를 kJ 로 바꿔도 비율은 그대로다.
\end{example}

% =====================================================================
\section{범주형 --- 명목과 순서}
\label{sec:attr-categorical}
% source: deck 3 slide 3

\begin{inbrief}
  \item[Why] 범주에 순서가 있으면 중앙값과 정렬이 허용되고, 없으면 금지되기
             때문이다.
  \item[When] 1단계에서 범주형으로 가른 뒤에 던지는 질문이다.
  \item[Where] 혈액형 $\cdot$ 지역 $\cdot$ 음식 종류는 명목, 학력 $\cdot$ 만족도
               $\cdot$ 학점은 순서다.
  \item[How] ``값들 사이에 자연스러운 서열이 있는가''를 묻는다. 나열 순서를
             뒤바꿔도 아무 문제가 없다면 명목이다.
\end{inbrief}

\begin{definition}[명목 척도와 순서 척도]
\label{def:nominal-ordinal}
\term{명목 척도}(Nominal Scale)는 ``같다/다르다''만 판정할 수 있는 척도다.
서열이 없으므로 나열 순서는 임의다.

\term{순서 척도}(Ordinal Scale)는 값들 사이에 서열이 있는 척도다. 단, 등급
사이의 \emph{간격이 같다는 보장은 없다}.
\end{definition}

순서 척도에서 가장 자주 나오는 실수가 이 간격 문제다. 만족도를
``매우 불만족 $=1$, 불만족 $=2$, 보통 $=3$, 만족 $=4$, 매우 만족 $=5$''로
숫자를 붙이면 평균을 계산할 수 있게 되지만, 그 평균 3.4는
``1$\to$2 의 차이''와 ``4$\to$5 의 차이''가 같다고 가정한 값이다. 그런 가정을
한 적은 없다.

\begin{pitfall}[숫자로 저장되었다고 수치형이 아니다]
우편번호, 학번, 제품 코드, 등번호는 전부 정수로 저장되지만 \term{명목
척도}(Nominal Scale)다. \texttt{pandas} 는 이들을 \texttt{int64} 로 읽고
\texttt{.mean()} 도 아무 불평 없이 계산해 준다. 오류가 나지 않는다는 것이
가장 위험한 지점이다 --- 유형은 저장 타입(dtype)이 아니라 \textbf{의미}로
정한다.
\end{pitfall}

% =====================================================================
\section{유형이 통계를 결정한다}
\label{sec:attr-stats}

\begin{inbrief}
  \item[Why] 속성 유형을 가르는 목적 자체가 어떤 요약값이 합법인지 정하는
             것이기 때문이다.
  \item[When] \Cref{ch:descriptive}의 요약값을 고르기 직전에 확인한다.
  \item[Where] 중심값 선택(\Cref{sec:center})과 그래프 선택
               (\Cref{ch:visualization})에 그대로 이어진다.
  \item[How] \Cref{tab:attr-ops} 의 행을 위에서 아래로 훑어 허용된 연산까지만
             쓴다.
\end{inbrief}

허용되는 연산은 \term{명목}$\to$\term{순서}$\to$\term{구간}$\to$\term{비율}
순으로 누적된다. 위 단계에서 허용된 것은 아래 단계에서도 전부 허용된다
(\Cref{tab:attr-ops}).

\begin{table}[ht]\centering
\caption{속성 유형별로 허용되는 연산과 대표 통계. 아래로 갈수록 허용 범위가
넓어진다.}
\label{tab:attr-ops}
\begin{tabularx}{\linewidth}{@{}l l l X@{}}
\toprule
척도 & 허용 연산 & 대표 통계 & 예시 \\
\midrule
명목(Nominal)  & $=,\ \neq$                    & 최빈값(Mode)
               & 혈액형 $\cdot$ 우편번호 $\cdot$ 음식 종류 \\
순서(Ordinal)  & $+\ <,\ >$                    & $+$ 중앙값(Median)
               & 학력 $\cdot$ 만족도 $\cdot$ 학점 \\
구간(Interval) & $+\ -$                        & $+$ 평균(Mean) $\cdot$ 표준편차
               & 섭씨온도 $\cdot$ 연도 $\cdot$ 지능지수 \\
비율(Ratio)    & $+\ \times,\ \div$            & $+$ 기하평균 $\cdot$ 변동계수
               & 칼로리 $\cdot$ 무게 $\cdot$ 금액 $\cdot$ 개수 \\
\bottomrule
\end{tabularx}
\end{table}

이 표는 \Cref{ch:descriptive} 전체의 전제다. 최빈값은 네 유형 모두에서 쓸 수
있고, 중앙값은 순서 척도부터, 평균과 표준편차는 구간 척도부터 허용된다. 즉
\Cref{sec:center}에서 ``평균이냐 중앙값이냐''를 고민하는 일 자체가 \textbf{구간
척도 이상일 때만} 성립하는 질문이다. 명목 척도에서는 그 질문이 아예 없다.

\begin{example}[과제 데이터의 여덟 변수]
\label{ex:hw1-types}
HW1 의 패스트푸드 영양정보 데이터
(\Cref{ch:hw1})는 여덟 개 변수를 가진다. 유형을 가르면 이렇게 된다.

\begin{center}
\begin{tabularx}{0.95\linewidth}{@{}l l X@{}}
\toprule
변수 & 유형 & 그렇게 가른 이유 \\
\midrule
\texttt{restaurant}   & 명목(Nominal) & 식당 이름. 순서가 없다 \\
\texttt{item}         & 명목(Nominal) & 메뉴 이름. 고유 식별자에 가깝다 \\
\texttt{type}         & 명목(Nominal) & 음식 종류. 우열이 없다 \\
\texttt{serving\_size}& 비율(Ratio)   & 0\,g 은 양이 없다는 뜻 \\
\texttt{calories}     & 비율(Ratio)   & 0\,kcal 은 열량이 없다는 뜻 \\
\texttt{fat}          & 비율(Ratio)   & 0\,g 은 지방이 없다는 뜻 \\
\texttt{sodium}       & 비율(Ratio)   & 0\,mg 은 나트륨이 없다는 뜻 \\
\texttt{sugars}       & 비율(Ratio)   & 0\,g 은 당류가 없다는 뜻 \\
\bottomrule
\end{tabularx}
\end{center}

수치형 다섯 개가 전부 비율 척도이므로 평균 $\cdot$ 표준편차 $\cdot$ 비율 계산이
모두 합법이다. 반면 명목 변수 셋에는 최빈값과 빈도만 쓸 수 있다 ---
\texttt{restaurant} 의 ``평균''은 존재하지 않는다.
\end{example}

\begin{aside}[이 데이터에 구간 척도가 없는 이유]
구간 척도는 사람이 기준점을 약속해서 만든 눈금(온도, 연도, 시험 점수)에서
주로 나온다. 물리적 양을 측정한 값은 대개 0 이 곧 ``없음''이므로 비율 척도가
된다. 영양정보처럼 측정값만 모인 데이터에 구간 척도가 없는 것은 우연이 아니다.
\end{aside}

% =====================================================================
\section{이산형과 연속형}
\label{sec:attr-discrete}
% source: deck 3 slide 2, 4

\begin{inbrief}
  \item[Why] 척도(명목$\cdot$순서$\cdot$구간$\cdot$비율)만으로는 값 사이에
             다른 값이 들어갈 수 있는지 알 수 없다. 이 축이 막대그래프
             (Bar Chart) 대 히스토그램(Histogram), 확률질량함수(PMF)
             대 확률밀도함수(PDF)의 선택을 가른다.
  \item[When] 척도를 정한 뒤, 그래프 종류나 확률 모형을 고르기 전에
              확인한다.
  \item[Where] 수치형과 범주형 모두에 적용되는 두 번째 축이다.
  \item[How] 두 값 사이에 중간값을 가진 다른 값이 실제로 존재할 수
             있는지 묻는다. 있으면 연속형, 셀 수 있게 끊겨 있으면
             이산형이다.
\end{inbrief}

값이 놓인 척도(\Cref{sec:attr-numeric}, \Cref{sec:attr-categorical})와
값이 끊겨 있는지는 \textbf{다른 질문}이다. 태어난 연도(구간 척도)와
키(비율 척도)는 둘 다 수치형이지만, 태어난 연도는 정수로만 끊겨 있고
키는 두 값 사이에 항상 더 세밀한 값이 있다. 이 차이를 놓치면 이산형
값을 연속형 그래프(히스토그램)로, 연속형 값을 이산형 그래프(막대그래프)로
그리는 실수가 나온다.

\begin{definition}[이산형과 연속형]
\label{def:discrete-continuous}
\term{이산형}(Discrete)은 값이 서로 떨어져 있어 셀 수 있는 변수다(예:
태어난 연도, 참/거짓). \term{연속형}(Continuous)은 실수 구간의 값을
갖는 변수다(예: 길이). 두 값 사이에 항상 더 세밀한 값을 넣을 수 있으면
연속형이다.
\end{definition}

척도(\Cref{tab:attr-ops})와 이산/연속은 서로 다른 축이므로, 실제 변수는
둘을 \textbf{함께} 답해야 완전히 분류된다. \Cref{tab:attr-discrete-2x2}가
이 두 축을 겹쳐 보여준다.

\begin{table}[ht]\centering
\caption{척도(행)와 이산/연속(열)을 겹친 분류. ``없음''은 그 조합이
존재하지 않는다는 뜻이다.}
\label{tab:attr-discrete-2x2}
\begin{tabularx}{\linewidth}{@{}l X X@{}}
\toprule
 & 이산형(Discrete) & 연속형(Continuous) \\
\midrule
수치형(Numeric) & 태어난 연도(구간), 고객 수(비율)
  & 길이 $\cdot$ 무게(비율), 섭씨온도 $\cdot$ 경도(구간) \\
범주형(Categorical) & 참/거짓(명목), 전공(명목), 만족도 1$\sim$5(순서)
  & 없음 \\
\bottomrule
\end{tabularx}
\end{table}

범주형 열이 연속형 칸에서 비는 이유는 한 문장으로 정리된다. 범주형
값은 정해진 기호의 집합에서 하나를 고르므로, 두 기호 사이에 들어갈
``중간 기호''가 없다. 그래서 범주형은 언제나 이산형이다.

\begin{example}[다섯 항목의 척도와 이산/연속을 함께 답하기]
\label{ex:attr-slide4}
다섯 개 항목에 척도와 이산/연속을 함께 매긴다(만족도 문항의 실제 모양은
1$\sim$5점 라디오 버튼이다). 척도 답은 슬라이드와 같고, 이산/연속은
이 책이 추가했다.

\begin{center}
\begin{tabularx}{\linewidth}{@{}l l l X@{}}
\toprule
항목 & 척도 & 이산/연속 & 근거 \\
\midrule
색(Color) & 명목(Nominal) & 이산형 & 색 이름 사이에 서열이 없다 \\
경도(Longitude) & 구간(Interval) & 연속형 & $0^\circ$는 Greenwich를
  고른 약속일 뿐이다 \\
강사 만족도 1$\sim$5 & 순서(Ordinal) & 이산형 & 등급 사이 간격이
  같다는 보장이 없다(\Cref{def:nominal-ordinal}) \\
무게(Weight) & 비율(Ratio) & 연속형 & 0\,kg은 무게가 없다는 뜻이다 \\
전공(College Major) & 명목(Nominal) & 이산형 & 전공 이름 사이에
  서열이 없다 \\
\bottomrule
\end{tabularx}
\end{center}

경도는 $180^\circ$와 $-180^\circ$가 같은 지점을 가리키는 원형 값이라,
두 경도의 ``차이''를 구할 때도 이 점을 조심해야 한다.
\end{example}

% =====================================================================
\section{표 데이터와 데이터 행렬}
\label{sec:data-matrix}
% source: deck 3 slide 5-8

\begin{inbrief}
  \item[Why] 열 하나씩 따로 보는 것만으로는 ``객체 하나가 여러 속성을
             동시에 가진다''는 사실을 다룰 수 없다. 행과 열을 한 번에
             다루는 틀이 필요하다.
  \item[When] 척도와 이산/연속을 다 가른 뒤, 표 전체를 계산에 넣기
              직전에.
  \item[Where] 행마다 하나의 대상(사람, 제품, 관측)을 담고, 열마다
               하나의 속성을 담는 모든 표.
  \item[How] 표를 $n \times d$ 행렬로 본다. $n$은 행(객체) 수, $d$는
             열(속성) 수다.
\end{inbrief}

열 하나씩 값을 따로 읽으면, 예를 들어 급여 열과 나이 열을 각각 볼 수는
있어도 ``이 사람의 급여와 나이''처럼 한 대상이 가진 여러 속성을 한꺼번에
계산에 넣을 방법이 없다. \term{표 데이터}(Tabular Data)는 행 하나가
대상 하나, 열 하나가 속성 하나를 담도록 정리해 이 문제를 없앤다. 이렇게
정리된 표는 $n \times d$ \term{데이터 행렬}(Data Matrix) $\mathbf{D} \in
\mathbb{R}^{n \times d}$로 쓸 수 있다. $n$은 행의 수(대상 수), $d$는
열의 수(속성 수)다.

행렬의 행 하나는
\[
  \mathbf{x}_i = (x_{i1}, x_{i2}, \dots, x_{id})^T \in \mathbb{R}^d
\]
로 쓴다. 대상 $i$번째의 속성 값 $d$개를 순서대로 늘어놓은 것이다.
행렬의 열 하나는
\[
  X_j = (x_{1j}, x_{2j}, \dots, x_{nj})^T \in \mathbb{R}^n
\]
로 쓴다. 속성 $j$번째의 값을 모든 대상에서 모은 것이다. \textbf{$\mathbf{x}_i$
(굵은 소문자, 행)와 $X_j$(대문자, 열)는 다른 것이다.} 첨자가 대상을
가리키면 행, 속성을 가리키면 열이다. 분야마다 행과 열을 부르는 이름이
다른데, \Cref{tab:attr-synonyms}가 그 대응을 정리한다.

\begin{table}[ht]\centering
\caption{데이터 행렬의 행과 열을 부르는 여러 이름. 같은 것을 분야마다
다르게 부른다.}
\label{tab:attr-synonyms}
\begin{tabularx}{\linewidth}{@{}l X@{}}
\toprule
대상 & 동의어 \\
\midrule
행($\mathbf{x}_i$) & 표본(Sample), 인스턴스(Instance), 레코드(Record),
  객체(Object), 점(Point), 특징 벡터(Feature Vector) \\
열($X_j$) & 속성(Attribute), 특징(Feature), 차원(Dimension),
  변수(Variable), 필드(Field) \\
\bottomrule
\end{tabularx}
\end{table}

\begin{example}[이름, 급여, 나이 표를 행렬로 보기]
\label{ex:attr-name-salary-age}
\Cref{tab:attr-name-salary-age}는 이름, 급여, 나이 세 열로 된 표다.

\begin{table}[ht]\centering
\caption{이름, 급여, 나이 표.}
\label{tab:attr-name-salary-age}
\begin{tabularx}{0.7\linewidth}{@{}l X X@{}}
\toprule
이름(Name) & 급여(Salary, \$) & 나이(Age, Years) \\
\midrule
Jane    & 90000 & 52 \\
John    & 85000 & 48 \\
Delilah & 75000 & 32 \\
Dave    & 90000 & 53 \\
Ellen   & 82000 & 44 \\
\bottomrule
\end{tabularx}
\end{table}

세는 순서는 한 줄씩이다.
\begin{align*}
  &\text{머리글 1행 제외} \to 5\text{행}
    && \text{Jane부터 Ellen까지 5명} \\
  &\text{열 3개}
    && \text{이름, 급여, 나이} \\
  &\text{이름(식별자) 열 제외} \to \text{속성 2개}
    && \text{급여, 나이만 속성이다} \\
  &\mathbf{D} \in \mathbb{R}^{5 \times 2}
    && \text{데이터 행렬 크기}
\end{align*}
\end{example}

\begin{pitfall}[식별자 열은 속성이 아니다]
\Cref{tab:attr-name-salary-age}의 이름(Name) 열처럼 대상을 가리키기만
하는 열은 속성으로 세지 않는다. 이 규칙을 놓치면 이 표와, 뒤에 나올
급여$\cdot$나이$\cdot$근속연수 표인 \Cref{tab:attr-3d-data}가 헷갈린다.
두 표 모두 ``5행 3열''이라 부를 수 있지만, 두 슬라이드가 세는 대상이
다르다.
이 표의 3열은 이름(식별자) 포함 3열이라 속성은 2개이고,
\Cref{tab:attr-3d-data}의 3열은 이름 열이 빠지고 근속연수가 더해진
3열이라 속성이 3개다. 규칙 하나로 정리하면 이렇다. \textbf{식별자는
속성이 아니다. 차원은 식별자를 뺀 열의 수다.} \Cref{ex:hw1-types}의
\texttt{item}도 고유값이 126행 중 107개뿐이라 완전한 식별자는 아니지만,
원칙은 같다. 대상 하나하나를 실제로 가리키는 것은 행 번호이지 그 열의
값이 아니다.
\end{pitfall}

% =====================================================================
\section{차원의 두 가지 뜻}
\label{sec:dimension}
% source: deck 3 slide 9-12

\begin{inbrief}
  \item[Why] 같은 낱말 ``차원''이 프로그래밍과 데이터 과학에서 다른
             것을 센다. 섞어 쓰면 ``속성이 100개인 데이터는 100차원
             배열이 필요하다''는 오해가 생긴다.
  \item[When] 데이터 행렬(\Cref{sec:data-matrix})의 크기를 말할 때마다.
  \item[Where] 저장 방식(배열)을 말할 때와 데이터가 놓인 공간을 말할
               때 모두에 쓰이는 말이라 헷갈린다.
  \item[How] ``지금 세는 것이 배열의 축인가, 속성의 개수인가''를 먼저
             구분한다.
\end{inbrief}

``차원''이라는 말 하나로 표의 축 개수와 속성 개수를 동시에 가리키면,
속성이 100개인 표를 ``100차원 배열''로 저장해야 한다는 오해가 생긴다.
표는 속성이 몇 개든 행과 열, 축 2개짜리 배열일 뿐이다. 이 오해를
없애려면 두 뜻을 다른 이름으로 나눠야 한다. \Cref{tab:attr-dimension-two-meanings}가
이 두 뜻을 나눈다. \term{배열 차원}(Array Dimension)은 값에 접근할
때 쓰는 인덱스(축)의 개수다. \term{공간 차원}(Space Dimension)은
데이터 행렬의 속성 수 $d$, 즉 각 점이 놓인 공간의 차원이다.

\begin{table}[ht]\centering
\caption{프로그래밍의 차원과 데이터 과학의 차원.}
\label{tab:attr-dimension-two-meanings}
\begin{tabularx}{\linewidth}{@{}l X@{}}
\toprule
뜻 & 세는 대상 \\
\midrule
배열 차원(Array Dimension) & 값에 접근하는 인덱스(축)의 수. 표(행,
  열)는 축이 2개이므로 항상 배열 차원 2다 \\
공간 차원(Space Dimension) & 속성의 수 $d$. 속성이 3개면 그 점들은
  3차원 공간에 놓인다 \\
\bottomrule
\end{tabularx}
\end{table}

\begin{pitfall}[슬라이드 9: 차원 = 셀의 수?]
슬라이드 값은 ``프로그래밍의 차원(Dimension) = 셀의 수''이고, 다시
계산하면 ``차원 = 인덱스(축)의 수''다. 표 하나는 셀이 아무리 많아도
축은 행과 열 둘뿐이라 배열 차원은 항상 2다. 슬라이드 11 자신도 아래에서
같은 표를 ``2차원 배열(2-dimensional array)''이라 부른다.
\end{pitfall}

\paragraph{2차원 예.} 이름/급여/나이 표인 \Cref{tab:attr-name-salary-age}에서
이름 열을 빼면 급여$\cdot$나이 두 속성만 남고, 각 사람은 2차원 평면의
점 하나가 된다(\Cref{fig:attr-2d-scatter}). Jane(52세, 90000)과
Dave(53세, 90000)는 나이만 한 살 다를 뿐이라 그림에서 거의 겹쳐
보인다.

\begin{figure}[ht]\centering
  \includegraphics[width=0.75\linewidth]{03_01_salary_age_2d.png}
  \caption{급여$\cdot$나이 2차원 산점도(Scatter Plot). 5명이 각각 점
  하나에 대응한다.}
  \label{fig:attr-2d-scatter}
\end{figure}

\paragraph{3차원 예.} 같은 다섯 사람에게 근속연수(Years in Service)를
더하면 \Cref{tab:attr-3d-data}가 된다. 이 표는 이름 없이 급여, 나이,
근속연수 세 열뿐이다.

\begin{table}[ht]\centering
\caption{급여, 나이, 근속연수 표.}
\label{tab:attr-3d-data}
\begin{tabularx}{0.6\linewidth}{@{}X X X@{}}
\toprule
급여(Salary, \$) & 나이(Age, Years) & 근속연수(Years in Service) \\
\midrule
90000 & 52 & 10 \\
85000 & 48 & 20 \\
75000 & 32 & 30 \\
90000 & 53 & 40 \\
82000 & 44 & 20 \\
\bottomrule
\end{tabularx}
\end{table}

\begin{align*}
  &5 \times 3 = 15 && \text{배열의 셀 수} \\
  &\text{인덱스(축) 2개(행, 열)} && \text{배열 차원 2} \\
  &\text{속성 3개(급여, 나이, 근속연수)} && \text{공간 차원 3}
\end{align*}

이름은 행렬에는 들어가지 않지만, 그림에서 점의 라벨로는 남는다
(\Cref{fig:attr-3d-scatter}).

\begin{figure}[ht]\centering
  \includegraphics[width=0.8\linewidth]{03_02_salary_age_service_3d.png}
  \caption{급여$\cdot$나이$\cdot$근속연수 3차원 산점도. 이름은 점의
  라벨로만 남는다.}
  \label{fig:attr-3d-scatter}
\end{figure}

\paragraph{속성이 100개라도 배열은 2차원이다.} 표 하나는 속성이 몇
개든 $n \times d$ 꼴이므로, $d = 100$이어도 여전히 축은 행과 열
둘뿐이라 배열 차원은 2다. 반면 그 점들이 놓인 공간 차원은 100이다.

\begin{pitfall}[슬라이드 12: 고차원 데이터 = 행이 많음?]
슬라이드 값은 ``고차원 데이터(High Dimension Data)는 행이 많은 것도
포함한다''이고, 다시 계산하면 ``고차원은 열(속성) 수 $d$가 클 때만
해당한다''다. 슬라이드 9 자신의 정의(차원 수 = 특징 수)를 그대로
적용하면, 행이 많은 것은 대상 수 $n$이 큰 대용량(Large-Scale) 데이터
이지 고차원 데이터가 아니다.
\end{pitfall}

\begin{pitfall}[슬라이드 9, 12: 마지막 열 머리글]
슬라이드 값은 표의 마지막 열 머리글 ``Feature $n$''이고, 다시
계산하면 ``Feature $d$''다. \Cref{sec:data-matrix}의 기호에서 $n$은
행(대상) 수이고 열(속성) 수를 가리키는 기호는 $d$이므로, 이 머리글은
기호가 뒤바뀐 것이다.
\end{pitfall}

% =====================================================================
\section{점, 행, 벡터}
\label{sec:vector}
% source: deck 3 slide 13-15

\begin{inbrief}
  \item[Why] 행을 그냥 ``숫자 목록''으로만 보면 두 대상 사이의 거리나
             평균 점 같은 기하 연산을 떠올릴 수 없다.
  \item[When] 거리, 평균, 방향처럼 기하학적으로 데이터를 다룰 때.
  \item[Where] 데이터 행렬의 모든 행. 물리의 벡터와 같은 방식으로
               계산한다.
  \item[How] 행 $\mathbf{x}_i$를 원점에서 그 점까지 뻗은 화살표로
             본다.
\end{inbrief}

행 $\mathbf{x}_i$를 숫자가 나열된 목록으로만 보면, 두 대상 사이의
거리나 다섯 대상의 평균 점 같은 기하 연산을 떠올릴 방법이 없다.
데이터 과학의 \term{벡터}(Vector)는 그 목록을 물리의 벡터와 똑같이,
표의 행 하나이자 공간의 점 하나로, 원점에서 그 점까지 뻗은 화살표로
본다.
\Cref{tab:attr-vectors-source}의 다섯 점(나이, 급여)을 화살표로 그리면
\Cref{fig:attr-vectors}가 된다.

\begin{table}[ht]\centering
\caption{나이, 급여 표(벡터 예).}
\label{tab:attr-vectors-source}
\begin{tabularx}{0.5\linewidth}{@{}X X@{}}
\toprule
나이(Age) & 급여(Salary) \\
\midrule
20 & 90000 \\
30 & 85000 \\
28 & 40000 \\
40 & 95000 \\
35 & 42000 \\
\bottomrule
\end{tabularx}
\end{table}

\begin{figure}[ht]\centering
  \includegraphics[width=0.75\linewidth]{03_03_vectors_origin.png}
  \caption{원점에서 각 점까지 뻗은 화살표. 점 하나가 벡터 하나다.}
  \label{fig:attr-vectors}
\end{figure}

\paragraph{붓꽃 예.} 이 데이터는 Fisher의 붓꽃 데이터(Iris Data)다.
자세한 출처는 \Cref{sec:qq-iris-source}에서 다룬다. 이 책은
\texttt{sklearn.datasets.load\_iris()}로 값만 읽는다.
\Cref{tab:attr-iris-rows}가 그중 10행을 옮긴다.

\begin{table}[ht]\centering
\caption{붓꽃 데이터 10행. $\mathbf{x}_1 \sim \mathbf{x}_8$은 setosa,
$\mathbf{x}_{149}, \mathbf{x}_{150}$은 virginica다.}
\label{tab:attr-iris-rows}
\begin{tabularx}{\linewidth}{@{}l X X X X l@{}}
\toprule
행 & 꽃받침 길이 & 꽃받침 너비 & 꽃잎 길이 & 꽃잎 너비 & 종(Class) \\
 & (Sepal Length) & (Sepal Width) & (Petal Length) & (Petal Width) & \\
\midrule
$\mathbf{x}_1$   & 5.1 & 3.5 & 1.4 & 0.2 & setosa \\
$\mathbf{x}_2$   & 4.9 & 3.0 & 1.4 & 0.2 & setosa \\
$\mathbf{x}_3$   & 4.7 & 3.2 & 1.3 & 0.2 & setosa \\
$\mathbf{x}_4$   & 4.6 & 3.1 & 1.5 & 0.2 & setosa \\
$\mathbf{x}_5$   & 5.0 & 3.6 & 1.4 & 0.2 & setosa \\
$\mathbf{x}_6$   & 5.4 & 3.9 & 1.7 & 0.4 & setosa \\
$\mathbf{x}_7$   & 4.6 & 3.4 & 1.4 & 0.3 & setosa \\
$\mathbf{x}_8$   & 5.0 & 3.4 & 1.5 & 0.2 & setosa \\
$\mathbf{x}_{149}$ & 6.2 & 3.4 & 5.4 & 2.3 & virginica \\
$\mathbf{x}_{150}$ & 5.9 & 3.0 & 5.1 & 1.8 & virginica \\
\bottomrule
\end{tabularx}
\end{table}

이 표 전체는 150행 5열이다. 수치 열은 꽃받침 길이$\cdot$너비, 꽃잎
길이$\cdot$너비 4개뿐이다. 종(Class) 열 $X_5$는 이름(setosa,
versicolor, virginica)을 담은 명목 범주형(\Cref{sec:attr-categorical})
이라, 수치 데이터 행렬 $\mathbf{D} \in \mathbb{R}^{150 \times 4}$에서는
뺀다.

\paragraph{평균 점(Mean Point).} 꽃받침 길이와 너비 두 열만 놓고 150개 점을 찍은
뒤 평균 점을 표시하면 \Cref{fig:attr-iris-mean}이 된다. 평균 점은
열마다 따로 평균을 낸 벡터 $(\bar X_1, \bar X_2)$다.

\begin{align*}
  &\bar X_1 = 5.8433 && \text{꽃받침 길이 150개의 평균} \\
  &\bar X_2 = 3.0573 && \text{꽃받침 너비 150개의 평균}
\end{align*}

\begin{figure}[ht]\centering
  \includegraphics[width=0.75\linewidth]{03_04_iris_points_mean.png}
  \caption{붓꽃 꽃받침 길이$\times$너비 150개 점과 평균 점
  $(5.8433,\ 3.0573)$.}
  \label{fig:attr-iris-mean}
\end{figure}

% =====================================================================
\section{무엇이든 행렬로}
\label{sec:everything-matrix}
% source: deck 3 slide 16-19

지금까지는 이미 표로 정리된 데이터를 다뤘다. 그런데 이미지, 텍스트,
네트워크처럼 표 모양이 아닌 데이터도 있다. 이런 데이터도 행과 열을
정해 같은 $n \times d$ 데이터 행렬 틀에 넣을 수 있다.

\paragraph{이미지.} 이미지 하나를 여러 영역 $r_1, r_2, \dots$로 나누고,
영역마다 RGB 색이나 질감(Texture) 같은 특징을 뽑으면, 영역이 행이
되고 특징이 열이 되는 행렬이 만들어진다(\Cref{fig:attr-image-matrix}).
슬라이드에 실린 원본 사진 두 장은 출처가 밝혀지지 않아 이 책에는
싣지 않는다.

\begin{figure}[ht]\centering
\begin{tikzpicture}[
  imgbox/.style = {draw=bookaccent, fill=bookaccentbg, rounded corners=2pt,
                    minimum width=20mm, minimum height=20mm, align=center,
                    font=\small},
  rbox/.style   = {draw=bookrule, fill=booksurface, rounded corners=2pt,
                    minimum width=16mm, minimum height=8mm, align=center,
                    font=\scriptsize},
  arrow/.style  = {-{Stealth[length=2.2mm]}, draw=bookmuted}]

  \node[imgbox] (img) at (0,0) {이미지\\(Image)};
  \node[rbox] (r1) at (3.0, 1.0) {영역 $r_1$};
  \node[rbox] (r2) at (3.0, -1.0) {영역 $r_2$};
  \draw[arrow] (img.east) -- (r1.west);
  \draw[arrow] (img.east) -- (r2.west);

  \node[align=left, font=\scriptsize] (tbl) at (7.6, 0) {%
    \begin{tabular}{@{}l l l@{}}
    \toprule
    영역 & RGB 색 & 질감 \\
    \midrule
    $r_1$ & \dots & \dots \\
    $r_2$ & \dots & \dots \\
    \bottomrule
    \end{tabular}};
  \draw[arrow] (r1.east) -- ++(0.5,0) |- ($(tbl.west)+(0,0.55)$);
  \draw[arrow] (r2.east) -- ++(0.5,0) |- ($(tbl.west)+(0,-0.55)$);
\end{tikzpicture}
\caption{이미지를 영역으로 나누고 영역마다 특징(RGB 색, 질감)을 뽑으면
영역이 행, 특징이 열인 데이터 행렬이 된다.}
\label{fig:attr-image-matrix}
\end{figure}

\paragraph{텍스트.} 단어 하나마다 그 단어와 함께 나타나는(동시출현,
Co-occurrence) 다른 단어들의 통계를 모아 벡터로 만든다. 슬라이드의
``King'' 벡터는 값이 50개다(앞 5개만 보이면 $0.50451$, $0.68607$,
$-0.59517$, $-0.022801$, $0.60046$, 뒤에 45개 더, 모두 50개). 슬라이드는 이
벡터 연산으로 king $-$ man $+$ woman $\approx$ queen이 성립한다고
보이지만, 슬라이드에는 king 벡터 하나만 실려 있어 이 책에서 다시
계산할 수는 없다. 이 기법에 붙는 모형 이름은 슬라이드에 나오지
않으므로 책이 임의로 붙이지 않는다.

\paragraph{네트워크.} 노드(Node) 여러 개와 그 사이를 잇는 간선(Edge)
으로 이뤄진 네트워크는 노드를 행과 열로 삼은 \term{인접 행렬}
(Adjacency Matrix)로 바꾼다. 두 노드 사이에 간선이 있으면 그 칸이 1,
없으면 0이다. \Cref{fig:attr-network}가 노드 5개짜리 네트워크를
인접 행렬로 바꾼 예다.

\begin{figure}[ht]\centering
\begin{tikzpicture}[
  netdot/.style = {draw=bookaccent, fill=bookaccentbg, circle,
                    minimum size=7mm, font=\small},
  eline/.style  = {draw=bookmuted, thick},
  arrow/.style  = {-{Stealth[length=2.5mm]}, draw=bookaccent, thick}]

  \node[netdot] (A) at (0, 1.3) {A};
  \node[netdot] (B) at (1.6, 1.3) {B};
  \node[netdot] (C) at (0, 0) {C};
  \node[netdot] (D) at (1.6, 0) {D};
  \node[netdot] (E) at (3.0, 0.65) {E};

  \draw[eline] (A) -- (B);
  \draw[eline] (A) -- (C);
  \draw[eline] (A) -- (D);
  \draw[eline] (B) -- (D);
  \draw[eline] (B) -- (E);
  \draw[eline] (C) -- (D);
  \draw[eline] (D) -- (E);
  \draw[eline] (D) to[out=300, in=330, looseness=8] (D);

  \draw[arrow] (4.0, 0.65) -- (5.0, 0.65);

  \node[font=\scriptsize] at (8.6, 0.65) {%
    $\begin{matrix}
       & A & B & C & D & E \\
      A & 0 & 1 & 1 & 1 & 0 \\
      B & 1 & 0 & 0 & 1 & 1 \\
      C & 1 & 0 & 0 & 1 & 0 \\
      D & 1 & 1 & 1 & 1 & 1 \\
      E & 0 & 1 & 0 & 1 & 0
     \end{matrix}$};
\end{tikzpicture}
\caption{노드 5개짜리 네트워크(왼쪽)를 인접 행렬(오른쪽)로 바꾼다.
D의 자기 고리(Self-loop)가 대각 성분 $A_{DD}=1$이 된다.}
\label{fig:attr-network}
\end{figure}

\begin{align*}
  &3+3+2+5+2 = 15 && \text{행 합(A, B, C, D, E 순)} \\
  &2 \times 7 + 1 = 15 && \text{간선 7개는 두 번, 자기 고리는 한 번
    센 값과 일치}
\end{align*}

인접 행렬은 대칭이다($\mathbf{A} = \mathbf{A}^T$). 간선에 방향이
없으므로 $i$행 $j$열 값과 $j$행 $i$열 값이 같기 때문이다.

% =====================================================================
\section{어떻게 적용하는가}
\label{sec:attr-apply}

HW1 의 패스트푸드 영양정보 데이터(\Cref{ch:hw1}, \Cref{ex:hw1-types})를
이 장의 틀로 세 번 센다.

\begin{align*}
  &126 \times 8 && \text{저장된 표 크기(행 126, 열 8)} \\
  &8 - 1 = 7 && \text{\texttt{item} 열을 빼면 속성 7개} \\
  &\mathbf{D} \in \mathbb{R}^{126 \times 5} && \text{수치 데이터
    행렬(비율 척도 5개 열)}
\end{align*}

\begin{pitfall}[\texttt{item} 열은 식별자가 아니다]
\texttt{item} 열은 식별자처럼 보이지만 126행에 고유값이 107개뿐이다.
같은 메뉴 이름이 여러 식당에 있기 때문이다(예: ``Cheeseburger''는
McDonald's와 Burger King 모두에 있다). 행을 하나로 가리키는 것은
행 번호이지 \texttt{item} 값이 아니다. 그래도 수치 데이터 행렬은
그대로 $\mathbf{D} \in \mathbb{R}^{126 \times 5}$다. \texttt{item}은
애초에 수치 열이 아니었기 때문이다.
\end{pitfall}

\paragraph{거리 하나로 보는 단위의 문제.} 2차원 예
(\Cref{fig:attr-2d-scatter})의 나이$\cdot$급여로 \term{유클리드
거리}(Euclidean Distance)를 재보면, 단위가 큰 급여가 거리를
압도한다.

\begin{align*}
  &\sqrt{1^2 + 0^2} = 1.0000 && \text{Jane–Dave 원래 척도 거리(나이
    1살, 급여 \$0 차이)} \\
  &\sqrt{4^2 + 5000^2} = 5000.0016 && \text{Jane–John 원래 척도 거리
    (나이 4살, 급여 \$5000 차이)}
\end{align*}

나이 차이 4는 급여 차이 \$5000 앞에서 거리에 거의 반영되지 않는다.
두 축을 각 축의 범위(나이 범위 21, 급여 범위 15000)로 나누는
\term{정규화}(Normalization)를 한 뒤 다시 재면, 어느 쪽이 더
가까운지(Jane–Dave가 Jane–John보다 가깝다)는 그대로다. 바뀌는 것은
나이가 거리에 차지하는 몫이다.

\begin{align*}
  &4 / 21 = 0.1905 && \text{Jane–John 정규화 나이 몫} \\
  &5000 / 15000 = 0.3333 && \text{Jane–John 정규화 급여 몫} \\
  &\sqrt{0.1905^2 + 0.3333^2} = 0.3839 && \text{Jane–John 정규화 거리} \\
  &1 / 21 = 0.0476 && \text{Jane–Dave 정규화 거리}
\end{align*}

정규화 전에는 급여 몫이 거리를 거의 전부 차지해(원래 척도 거리
5000.0016 중 급여가 5000) 나이는 있으나 마나였다. 정규화 뒤에는
나이 몫 0.1905와 급여 몫 0.3333이 같은 척도(0에서 1)에서 거리에 들어간다.
Jane–Dave가 Jane–John보다 가깝다는 순서는 정규화 전후로 바뀌지
않는다(0.0476 < 0.3839, 1.0000 < 5000.0016). 정규화를 언제,
어떻게 하는지는 \Cref{sec:normalization}에서 다룬다.

% =====================================================================
\section{앞으로}
\label{sec:attr-future}

\begin{itemize}
  \item \textbf{고차원(High Dimension)과 차원의 저주(Curse of
        Dimensionality).} 속성 수 $d$가 커지면 점들 사이의 거리가
        비슷해지고 계산량도 늘어난다. 이 문제와 그 처방인 차원 축소
        (Dimensionality Reduction)는 이 덱에 없다. 강의계획서(덱 1)
        슬라이드 4의 주제 목록에만 이름이 올라 있다.
  \item \textbf{이산화(Discretization).} 연속형 속성을 이산형 속성으로 바꾸는 다리는
        \Cref{sec:discretization}에서 다룬다. \Cref{tab:attr-discrete-2x2}
        의 열을 오른쪽에서 왼쪽으로 옮기는 것과 같다.
  \item \Cref{ch:hypothesis}의 가설검정은 변수가 어떤 척도인지 이
        장에서 정한 판단을 그대로 전제로 삼는다.
\end{itemize}
